\chapter{Introducción}
\label{chap:introduccion}

En este capítulo se presenta el contexto general del proyecto, junto con la motivación que ha llevado a su desarrollo. Para ello, primero se introducen los antecedentes relacionados con el estudio de la cobertura en redes celulares y la necesidad de herramientas de simulación visuales. A continuación, se describen los objetivos principales y específicos del proyecto y la metodología utilizada en su desarrollo. Por último, se muestra la organización de la memoria, indicando el contenido de cada capítulo.

\section{Motivación y antecedentes}

En los últimos años, las redes celulares han evolucionado notablemente debido al gran aumento de dispositivos conectados y a la necesidad de ofrecer servicios con mayores prestaciones. El \textit{streaming} de vídeo, las comunicaciones vehiculares o las comunicaciones críticas requieren redes capaces de proporcionar buena cobertura, baja latencia y una conexión estable. Por todo esto, el estudio de la propagación es fundamental para entender cómo se comporta la señal y cómo varía la cobertura en distintos entornos.

La cobertura de una red celular puede analizarse a partir de métricas como la potencia recibida, la relación señal-ruido (SNR) o la relación señal-interferencia más ruido (SINR). Estas métricas permiten determinar si la señal llega con un nivel suficiente y si puede distinguirse correctamente frente al ruido y las interferencias. En este proyecto no se modelan interferencias entre varios transmisores, por lo que el análisis se basa en la potencia recibida y en la SNR calculada en distintos puntos del escenario.

El análisis de la cobertura no depende únicamente de la distancia entre el transmisor y el receptor; también influyen otros factores, como la frecuencia utilizada, la potencia de transmisión, el patrón de radiación de las antenas o la presencia de obstáculos en el entorno. En escenarios tridimensionales, estos factores pueden provocar variaciones importantes en la señal recibida entre puntos cercanos, especialmente cuando existen edificios u otros elementos que provocan pérdidas adicionales.

Este tipo de estudios se realizan mediante modelos matemáticos, simulaciones numéricas o herramientas de representación bidimensional. Estas representaciones permiten analizar el comportamiento general de la señal, pero pueden dificultar la interpretación espacial de los resultados y omitir información relevante del escenario.

Uno de los aspectos más importantes que pueden perderse en una representación bidimensional es la influencia de la altura. En un entorno real, la señal no se propaga únicamente sobre un plano, sino que lo hace en tres dimensiones. La posición vertical del receptor, la altura de la antena transmisora, la orientación del patrón de radiación y la presencia de edificios pueden modificar de forma considerable la potencia recibida.

Además, la ganancia de una antena no tiene por qué ser igual en todas las direcciones. Por tanto, para analizar correctamente la cobertura, es necesario considerar tanto la distancia al receptor como la dirección de propagación.

A partir de estas necesidades surge la motivación principal de este proyecto: desarrollar un simulador tridimensional que permita analizar visualmente la cobertura en un entorno urbano. La representación 3D facilita relacionar los valores calculados con la posición real de los elementos del escenario y estudiar la influencia de la altura, los obstáculos y el patrón de radiación.

Para ello, el proyecto combina Unity con un simulador existente desarrollado en Python. Unity se utiliza para construir el escenario tridimensional, representar los mapas de calor y visualizar los receptores móviles, mientras que Python realiza los cálculos de propagación y de las métricas del enlace. De esta forma, la herramienta permite observar cómo varían la potencia recibida y la SNR en distintos puntos y recorridos del escenario.


\section{Objetivos del proyecto}

El objetivo principal de este trabajo es desarrollar un simulador tridimensional en Unity que permita analizar la cobertura de una red celular en un escenario urbano. La herramienta busca representar de forma visual cómo varía la señal en el entorno, teniendo en cuenta la posición del transmisor, el patrón de radiación de la antena, el canal de propagación, los obstáculos y la posición de los receptores.

Para alcanzar este objetivo principal, se plantean los siguientes objetivos específicos:

\begin{itemize}
    \item Construir un escenario urbano en tres dimensiones, incluyendo elementos como edificios, vehículos y peatones.
    \item Modelar el comportamiento de una antena transmisora con un modelo omnidireccional y posteriormente con un patrón de radiación específico.
    \item Reconstruir una aproximación tridimensional del patrón de radiación a partir de los cortes horizontal y vertical de una antena.
    \item Integrar Unity con un simulador ya existente desarrollado en Python para intercambiar los datos de la simulación y obtener métricas como la potencia recibida y la SNR en distintos puntos del escenario.
    \item Representar la cobertura en distintos puntos del escenario mediante un mapa de calor en 3D.
    \item Incorporar receptores móviles para analizar la variación de la cobertura durante su desplazamiento.
    \item Exportar los resultados obtenidos en formato CSV para facilitar su análisis posterior.
\end{itemize}

Junto a los objetivos específicos anteriores, también se plantearon una serie de mejoras opcionales para completar el simulador. Estas funcionalidades no eran imprescindibles para cumplir el objetivo principal, pero aportan al resultado final del proyecto:

\begin{itemize}
    \item Adaptar el escenario urbano genérico a una representación tridimensional del Centro Universitario de Mérida.
    \item Mejorar la visualización del mapa de calor añadiendo una vista bidimensional por capas de altura.
    \item Generar gráficas automáticamente a partir de los archivos CSV exportados.
    \item Añadir modos de rendimiento para mejorar la fluidez de la aplicación en dispositivos con prestaciones reducidas.
\end{itemize}

Con estos objetivos, se pretende crear una herramienta que no solo permita realizar cálculos de cobertura, sino también interpretar los resultados de una forma visual y accesible.

\section{Metodología de trabajo}

La metodología seguida en el proyecto fue incremental, dividiendo el trabajo en distintas fases. El desarrollo no se realizó como una única implementación completa desde el principio, sino como una evolución progresiva desde una primera versión sencilla hasta una más completa.

En primer lugar, se realizó un estudio inicial del problema y de los conceptos necesarios para entender mejor el análisis de cobertura en redes celulares, la propagación de la señal, los modelos de canal, los patrones de radiación de antenas y las métricas utilizadas para evaluar la cobertura.

Tras esta fase inicial, se analizó el simulador existente desarrollado en Python, con el objetivo de entender qué datos necesitaba como entrada y qué resultados podía devolver. A partir de ahí, se diseñó un sistema de comunicación entre Unity y Python para enviar desde Unity los datos de la simulación y realizar en Python los cálculos relacionados con la propagación y las métricas de cobertura.

El desarrollo se llevó a cabo por fases. Primero se construyó el escenario urbano genérico tridimensional y se implementó una primera versión de la simulación con una antena omnidireccional. Posteriormente, se añadió a la antena un patrón de radiación específico, reconstruyendo una aproximación tridimensional a partir de sus cortes horizontal y vertical. A continuación, se implementó la visualización de los resultados mediante mapas de calor, permitiendo representar de forma gráfica la cobertura obtenida en el escenario.

Después, se añadieron los receptores móviles y se calcularon sus métricas durante el desplazamiento por el escenario. Esta fase permitió analizar la evolución de la potencia recibida y la SNR para distintos recorridos. Finalmente, se incorporó la exportación de resultados en formato CSV, la generación automática de gráficas y diferentes mejoras visuales y de rendimiento.

Durante todo el proceso, se realizaron pruebas para comprobar el funcionamiento de cada fase antes de integrarla con el resto del sistema. De esta forma, se pudo validar paso a paso la comunicación entre Unity y Python, la reconstrucción del patrón de radiación, el cálculo de métricas y la representación visual de la cobertura. Esta forma de trabajo permitió avanzar de manera ordenada y detectar errores de forma más sencilla durante el desarrollo.

\section{Organización de la memoria}

La memoria está organizada en siete capítulos. Tras este primer capítulo introductorio, en el que se han presentado la motivación, los objetivos y la metodología seguida durante el proyecto, el resto del documento se estructura de la siguiente forma.

El segundo capítulo muestra el análisis y la descripción del proyecto. En él se plantea el problema que se pretende resolver, se definen los requisitos del sistema y se describen las tecnologías empleadas. También se incluye la planificación del trabajo y una valoración de la viabilidad del mismo.

El tercer capítulo corresponde al estado del arte. En él se explican los conceptos teóricos necesarios para comprender el proyecto, como las redes celulares, los modelos de propagación, el modelado de antenas, las comunicaciones vehiculares, la maximización de cobertura y las herramientas de simulación tridimensional.

El cuarto capítulo describe el diseño del simulador. En él se explica la arquitectura general de la aplicación, el flujo de ejecución, el diseño de la interfaz de configuración, la escena de simulación, la integración entre Unity y Python y el modelo de intercambio de información entre ambos entornos.

El quinto capítulo se centra en el desarrollo e implementación del simulador. Se detallan los principales elementos implementados, como la reconstrucción del patrón de radiación, el cálculo de ganancias, la generación de mapas de calor, los receptores móviles, la exportación de resultados y las mejoras de rendimiento.

El sexto capítulo presenta las pruebas realizadas y los resultados obtenidos. En él se analiza el funcionamiento del simulador mediante distintos escenarios y configuraciones, estudiando la cobertura obtenida, el comportamiento del patrón de radiación y la evolución de los receptores móviles.

Por último, el séptimo capítulo recoge las conclusiones del trabajo y plantea posibles líneas de mejora para trabajos futuros.